% ============================================================
%  Projeto Final -- Estrategia Quantitativa Completa
%  Nucleo Quant · Liga de Mercado Financeiro · UFU · 2026.1
% ============================================================
\documentclass[11pt,a4paper]{article}

% ─── Pacotes ────────────────────────────────────────────────────────────────
\usepackage[utf8]{inputenc}
\usepackage[T1]{fontenc}
\usepackage[top=3.5cm, bottom=2.5cm, left=2.5cm, right=2.5cm]{geometry}
\usepackage{xcolor}
\usepackage{titlesec}
\usepackage{enumitem}
\usepackage{amsmath, amssymb}
\usepackage{graphicx}
\usepackage{fancyhdr}
\usepackage{tcolorbox}
\usepackage{tabularx}
\usepackage{booktabs}
\usepackage{array}
\usepackage{float}
\usepackage{hyperref}

\tcbuselibrary{skins, breakable}

% ─── Cores ──────────────────────────────────────────────────────────────────
\definecolor{quantpurple}{HTML}{7B2FBE}
\definecolor{quantpurplelight}{HTML}{F0E8FF}
\definecolor{quantdark}{HTML}{1A1A1A}
\definecolor{quantgray}{HTML}{6B7280}
\definecolor{quantlightgray}{HTML}{F3F4F6}
\definecolor{quantgreen}{HTML}{16A34A}
\definecolor{quantred}{HTML}{DC2626}
\definecolor{quantamber}{HTML}{D97706}
\definecolor{quantblue}{HTML}{1D4ED8}
\definecolor{quantbluelight}{HTML}{EFF6FF}

% ─── Fonte sans-serif ────────────────────────────────────────────────────────
\renewcommand{\familydefault}{\sfdefault}

% ─── Hyperlinks ─────────────────────────────────────────────────────────────
\hypersetup{hidelinks}

% ─── Espacamento ─────────────────────────────────────────────────────────────
\setlength{\parskip}{4pt}
\setlength{\parindent}{0pt}

% ─── Cabecalho e rodape ──────────────────────────────────────────────────────
\setlength{\headheight}{2cm}
\setlength{\headsep}{1cm}

\pagestyle{fancy}
\fancyhf{}
\renewcommand{\headrulewidth}{0pt}
\fancyhead[L]{%
  \includegraphics[height=1.40cm]{logoSimples.png}%
}
\fancyhead[R]{%
  \raisebox{0.5cm}{\textcolor{quantgray}{\footnotesize
    N\'ucleo Quant
    \,$\cdot$\, Liga de Mercado Financeiro
    \,$\cdot$\, UFU \,$\cdot$\, 2026.1}}%
}
\fancyfoot[C]{\textcolor{quantgray}{\footnotesize\thepage}}
\fancyfoot[R]{%
  \textcolor{quantgray}{\footnotesize
    Projeto Final \,$\cdot$\, Estrat\'egia Quantitativa Completa}%
}

% ─── Titulos de secao ────────────────────────────────────────────────────────
\titleformat{\section}
  {\large\bfseries\sffamily\color{quantpurple}}
  {}{0em}{}
  [\vspace{2pt}\textcolor{quantpurple}{\rule{\textwidth}{1.2pt}}\vspace{4pt}]

\titlespacing{\section}{0pt}{18pt}{6pt}

% ─── Listas ──────────────────────────────────────────────────────────────────
\setlist[itemize]{
  label=\textcolor{quantpurple}{\small$\blacktriangleright$},
  leftmargin=1.8em,
  itemsep=3pt,
  parsep=0pt
}
\setlist[enumerate]{
  label=\textcolor{quantpurple}{\textbf{\arabic*.}},
  leftmargin=1.8em,
  itemsep=5pt,
  parsep=0pt
}

% ─── Codigo inline ───────────────────────────────────────────────────────────
\newcommand{\code}[1]{%
  \colorbox{quantlightgray}{\texttt{\small\color{quantdark}#1}}%
}

% ─── Caixa de destaque (roxo claro) ──────────────────────────────────────────
\newtcolorbox{highlightbox}{
  enhanced,
  colback=quantpurplelight,
  colframe=quantpurple,
  boxrule=1pt,
  arc=4pt,
  left=10pt, right=10pt, top=7pt, bottom=7pt,
  fontupper=\small\sffamily, coltext=quantdark,
}

% ─── Caixa de secao ──────────────────────────────────────────────────────────
\newtcolorbox{stepbox}[1]{
  enhanced,
  colback=white,
  colframe=quantpurple,
  boxrule=0.8pt,
  arc=4pt,
  title={\textbf{\textcolor{white}{\small\sffamily #1}}},
  attach boxed title to top left={yshift=-1mm, xshift=10pt},
  boxed title style={
    colback=quantpurple, colframe=quantpurple,
    arc=3pt, boxrule=0pt,
    left=6pt, right=6pt, top=2pt, bottom=2pt
  },
  left=10pt, right=10pt, top=14pt, bottom=8pt,
  fontupper=\small\sffamily, coltext=quantdark,
  breakable,
}

% ─── Caixa de atencao (ambar) ────────────────────────────────────────────────
\newtcolorbox{warningbox}{
  enhanced,
  colback=white,
  colframe=quantamber,
  boxrule=0.8pt,
  arc=4pt,
  title={\textbf{\textcolor{white}{\small\sffamily Aten\c{c}\~ao}}},
  attach boxed title to top left={yshift=-1mm, xshift=10pt},
  boxed title style={
    colback=quantamber, colframe=quantamber,
    arc=3pt, boxrule=0pt,
    left=6pt, right=6pt, top=2pt, bottom=2pt
  },
  left=10pt, right=10pt, top=14pt, bottom=8pt,
  fontupper=\small\sffamily, coltext=quantdark,
}

% ════════════════════════════════════════════════════════════════════════════
\begin{document}
% ════════════════════════════════════════════════════════════════════════════

\thispagestyle{fancy}

% ─── Titulo ──────────────────────────────────────────────────────────────────
\begin{center}
  \includegraphics[height=3cm]{logoLiga.png}\\[1em]
  {\fontsize{22}{26}\selectfont\bfseries\sffamily PROJETO FINAL}\\[0.4em]
  {\large\sffamily\textcolor{quantgray}{Estrat\'egia Quantitativa Completa}}
\end{center}

\vspace{0.5em}
\textcolor{quantpurple}{\rule{\textwidth}{2pt}}
\vspace{1em}

% ─── Informacoes ─────────────────────────────────────────────────────────────
\begin{tabularx}{\textwidth}{@{}lX@{}}
  \toprule
  \textcolor{quantpurple}{\textbf{Bloco}} &
  \textcolor{quantpurple}{\textbf{Formato}} \\
  \midrule
  Encerramento do ciclo 2026.1 & Grupos de 2 a 4 pessoas \\
  \bottomrule
\end{tabularx}

\vspace{1em}

% ─── Contexto ────────────────────────────────────────────────────────────────
\section{Contexto}

O Projeto Final \'e a s\'intese de tudo que foi aprendido no ciclo.
Voc\^e escolhe um ativo ou universo de ativos, formula uma hip\'otese
sobre o mercado, constr\'oi um modelo para explor\'a-la, testa rigorosamente
se ela funciona e apresenta os resultados com honestidade.

\begin{highlightbox}
  \textbf{O fluxo obrigat\'orio do projeto:}\\[6pt]
  \textbf{Tese} \textcolor{quantgray}{$\longrightarrow$}
  \textbf{Modelo} \textcolor{quantgray}{$\longrightarrow$}
  \textbf{Backtest} \textcolor{quantgray}{$\longrightarrow$}
  \textbf{Resultado}
  
\end{highlightbox}

\newpage

% ─── Estrutura do Projeto ────────────────────────────────────────────────────
\section{Estrutura do Projeto}

% ── PARTE 1 ──────────────────────────────────────────────────────────────────
\begin{stepbox}{Parte 1 \quad Tese}

A tese \'e a hip\'otese que motiva toda a estrat\'egia.
\'E uma afirma\c{c}\~ao sobre o mercado que voc\^e acredita ser verdadeira
e que, se verdadeira, gera oportunidade de lucro.

\vspace{0.4em}
A tese deve aparecer nos slides da apresenta\c{c}\~ao,
respondendo claramente a estas perguntas:

\begin{enumerate}
  \item \textbf{Qual \'e a ideia central?}\\
        Exemplos: \emph{``A\c{c}\~oes com momentum positivo de 12 meses
        tendem a continuar subindo no m\^es seguinte''},
        ou \emph{``Alta volatilidade hoje prediz alta volatilidade amanh\~a''}.

  \item \textbf{Por que essa inefici\^encia existiria?}\\
        Qual \'e o mecanismo econ\^omico ou comportamental que a sustenta?
        Exemplos: viés de ancoragem, reação lenta a notícias,
        prêmio de liquidez, etc.

  \item \textbf{Qual \'e o ativo ou universo de ativos?}\\
        Seja específico: ações do Ibovespa, um par de moedas,
        futuros de commodities, etc.

  \item \textbf{Qual \'e o horizonte de opera\c{c}\~ao?}\\
        Di\'ario, semanal, mensal?
\end{enumerate}
\end{stepbox}

\vspace{1em}

% ── PARTE 2 ──────────────────────────────────────────────────────────────────
\begin{stepbox}{Parte 2 \quad Modelo}

O modelo \'e a implementa\c{c}\~ao da tese em c\'odigo.
Ele define \textbf{quando comprar}, \textbf{quando vender}
e \textbf{quanto alocar}. Voc\^e tem liberdade de escolha --
use o que fizer mais sentido para sua tese:

\vspace{0.4em}
\begin{tabularx}{\textwidth}{@{}lX@{}}
  \toprule
  \textcolor{quantpurple}{\textbf{Tipo}} &
  \textcolor{quantpurple}{\textbf{Exemplos}} \\
  \midrule
  Estat\'istico &
  M\'edias m\'oveis, bandas de Bollinger, reversão à média \\[3pt]
  S\'eries temporais &
  ARIMA para previs\~ao, GARCH para volatilidade \\[3pt]
  Machine Learning &
  Random Forest, regressão logística com features de mercado \\[3pt]
  Baseado em risco &
  Portf\'olio de m\'inima vari\^ancia, paridade de risco \\
  \bottomrule
\end{tabularx}

\vspace{0.8em}
Independentemente do tipo, o modelo deve estar implementado em Python,
documentado com coment\'arios e gerar os \textbf{sinais de entrada e sa\'ida}.
\end{stepbox}

\vspace{1em}

% ── PARTE 3 ──────────────────────────────────────────────────────────────────
\begin{stepbox}{Parte 3 \quad Backtest}

O backtest simula a estrat\'egia no passado para avaliar se teria
funcionado. \'E a parte mais cr\'itica do projeto e a mais f\'acil
de fazer errado.

\begin{enumerate}
  \item \textbf{Per\'iodo m\'inimo:} 5 anos de dados.
        
  \item \textbf{Divis\~ao temporal:}
        \begin{itemize}
          \item \textbf{In-sample} (treino): per\'iodo usado para calibrar
                o modelo ou escolher par\^ametros
          \item \textbf{Out-of-sample} (teste): per\'iodo reservado
                exclusivamente para avalia\c{c}\~ao final.
                \textbf{Nunca use o per\'iodo de teste para ajustar par\^ametros}
        \end{itemize}

  \item \textbf{M\'etricas obrigat\'orias} (você pode e deve usar mais, dependendo do que sua estratégia precisar para se provar):
        \begin{itemize}
          \item \textbf{Retorno total} e \textbf{retorno anualizado}
          \item \textbf{Volatilidade anualizada}
          \item \textbf{\'Indice de Sharpe}
          \item \textbf{Maximum Drawdown} e dura\c{c}\~ao da recupera\c{c}\~ao
          \item \textbf{N\'umero de opera\c{c}\~oes} e \textbf{taxa de acerto}
                (porcentagem de trades lucrativos)
        \end{itemize}

  \item \textbf{Benchmark:} compare sempre com o \'indice que mais
        se adequa \`a sua estrat\'egia -- por exemplo, o \code{BOVA11.SA}
        para a\c{c}\~oes brasileiras, o \code{S\&P 500} para a\c{c}\~oes
        americanas, e assim por diante. A estrat\'egia bateu o \'indice
        em retorno? E em retorno ajustado ao risco (Sharpe)?

  \item \textbf{Custos de transa\c{c}\~ao:} considere que toda movimenta\c{c}\~ao
        envolve o pagamento de comiss\~ao (corretagem e spread).
        N\~ao incluir esses custos traria resultados irreais --
        estrat\'egias com muitas opera\c{c}\~oes s\~ao fortemente impactadas
        por custos ignorados.
\end{enumerate}

\begin{warningbox}
  \textbf{Vieses que invalidam um backtest:}
  \begin{itemize}
    \item \textbf{Look-ahead bias:} usar informa\c{c}\~ao futura no sinal.
          Ex.: calcular a m\'edia do dia \emph{t} usando o pre\c{c}o de
          fechamento do pr\'oprio dia \emph{t} para tomar a decis\~ao de
          comprar no fechamento de \emph{t}
    \item \textbf{Overfitting:} ajustar par\^ametros at\'e o backtest parecer
          bom -- e o modelo n\~ao generalizar para novos dados
    \item \textbf{Ignorar custos} 
  \end{itemize}
\end{warningbox}
\end{stepbox}

\vspace{1em}

% ── PARTE 4 ──────────────────────────────────────────────────────────────────
\begin{stepbox}{Parte 4 \quad Resultado e Conclus\~ao}

A conclus\~ao \'e onde voc\^e \textbf{interpreta} o que os n\'umeros
significam -- n\~ao apenas os lista. Deve responder a:

\begin{enumerate}
  \item A tese foi confirmada ou refutada pelos dados?

  \item Os resultados s\~ao economicamente significativos?
        Um Sharpe de 0,3 out-of-sample, sem custos, provavelmente
        n\~ao \'e explor\'avel na pr\'atica.

  \item Quais s\~ao as \textbf{limita\c{c}\~oes} da sua an\'alise?
        Exemplos: universo pequeno de ativos, par\^ametros sens\'iveis,
        per\'iodo atípico, liquidez n\~ao modelada.

  \item O que voc\^e faria diferente com mais tempo?
        Pr\'oximos passos concretos: outros ativos, outras features,
        walk-forward analysis, inclus\~ao de custos reais de corretagem.
\end{enumerate}
\end{stepbox}

\newpage

% ─── Defesa Presencial ───────────────────────────────────────────────────────
\section{Defesa Presencial}

A defesa presencial \'e o principal componente do projeto. Cada grupo
far\'a uma \textbf{apresenta\c{c}\~ao de 10 minutos com slides}, seguida de perguntas ao final.

\vspace{0.5em}
A defesa deve cobrir os quatro elementos do fluxo:

\vspace{0.4em}
\begin{tabularx}{\textwidth}{@{}lX@{}}
  \toprule
  \textcolor{quantpurple}{\textbf{Elemento}} &
  \textcolor{quantpurple}{\textbf{O que apresentar}} \\
  \midrule
  \textbf{Tese}      & A hip\'otese. Por que faz sentido? \\[4pt]
  \textbf{Modelo}    & Como foi implementado. Quais foram as decis\~oes t\'ecnicas? \\[4pt]
  \textbf{Backtest}  & Os gr\'aficos de equity curve e as m\'etricas principais \\[4pt]
  \textbf{Resultado} & Funcionou? Por qu\^e sim ou por qu\^e n\~ao? \\
  \bottomrule
\end{tabularx}

\vspace{1em}

% ─── Como Entregar ───────────────────────────────────────────────────────────
\section{Como Entregar o Código}

\begin{stepbox}{Estrutura do Reposit\'orio}

Crie um reposit\'orio \textbf{p\'ublico} no GitHub com o c\'odigo do projeto e envie no WhatsApp do Gerente do Quant até 23:59 da véspera da apresentação presencial.
Voc\^e pode organizar a estrat\'egia em um \'unico script ou em v\'arios
arquivos -- como fizer mais sentido. Os \'unicos arquivos
\textbf{obrigat\'orios} s\~ao:

\vspace{0.6em}
\begin{tabularx}{\textwidth}{@{}lX@{}}
  \toprule
  \textcolor{quantpurple}{\textbf{Arquivo}} &
  \textcolor{quantpurple}{\textbf{Conte\'udo}} \\
  \midrule
  \code{README.md}        & Descri\c{c}\~ao da estrat\'egia e instru\c{c}\~oes de execu\c{c}\~ao \\[3pt]
  \code{requirements.txt} & Bibliotecas com vers\~oes exatas \\
  \bottomrule
\end{tabularx}

\vspace{0.8em}

\begin{enumerate}
  \item \textbf{Escreva o README.md} com:
        \begin{itemize}
          \item Nome(s) e curso(s)
          \item A tese em 2--3 frases
          \item Como rodar: \code{pip install -r requirements.txt}
                e como executar o(s) script(s)
          \item A principal conclus\~ao do projeto
        \end{itemize}

  \item \textbf{Gere o requirements.txt:}\quad
        \code{pip freeze > requirements.txt}
\end{enumerate}
\end{stepbox}

% ─── Rodape ──────────────────────────────────────────────────────────────────
\vfill


\end{document}